
\emph{Paso 2. Variación del desplazamiento.} El Teorema~\ref{teo:insercion-torres} y el
Corolario~\ref{cor:dos-casillas}, aplicados al menor de $d$, dan en el soporte común
\begin{equation}\label{eq:variaciones-base}
 \delta p=O(ph),\qquad \delta G^0_{ij}=O(h^2),\qquad \delta\hat c=O(h^2),\qquad
 \delta\tau=O(ph),\qquad \delta\mu_i=O(h^2),
\end{equation}
la última por suma en una línea que pierde a lo sumo un sitio. De $\zeta=O(h^2)$ y
$s_x=-\mathsf G(x,e_0)+a_x^\top\zeta$ resulta primero $\delta s=O(h^2)$. En~\eqref{eq:masas-hijo} el
sumando de $B_i$ es $p_xs_x^2$ por una función con derivadas acotadas en $(p_x,s_x)$ cerca del origen.
En los índices comunes cambia en $O(p_xh^3)$, y el sitio perdido aporta $O(h^3)$; así $\delta
B_i=O(h^3)$. En la estrella $F^*$, la fórmula exacta~\eqref{eq:rho-hijo} da $\delta\rho_i=O(h^3)$. De los
dos sistemas lineales~\eqref{eq:masas-hijo}, para $N$ y para $N_d$, se restringe el mayor a los índices
comunes. Las sumas absolutas por fila de la diferencia de las dos matrices, incluido el acoplamiento
con los dos índices retirados, son $O(h)$: esto se sigue de $\delta\tau=O(ph)$, del sitio perdido y del
cambio de la proyección del calibre. Multiplicando por la solución, que es $O(h^2)$, y por la inversa,
uniformemente acotada en la norma del supremo, obtenemos
\[
 \delta\zeta=O(h^3),\qquad\text{y del mismo modo}\qquad \delta z=O(h^3),\qquad \delta u=O(h^2).
\]

\emph{Paso 3. Variación del resto.} Por el Paso~1 basta comparar las once expresiones en $N$ y en
$N_d$. Los términos interiores de los puntos 4, 6 y 7 tienen mayorantes $Cph^3$, $Cp^2h^2$ y $Cp^3h$
por sitio; cada uno suma $O(h^2)$, y $O(h^3)$ en las dos líneas perdidas. En los índices comunes el
telescopio de cada producto gana un factor $h$ en el mayorante, por~\eqref{eq:variaciones-base},
$\delta\alpha=O(h^2)$, $\delta z=O(h^3)$ y $\delta u=O(h^2)$; los factores escalares de orden $p^m$,
$m\ge1$, son analíticos con derivada $O(p^{m-1})$. En la estrella se pierden a lo sumo dos términos de
tamaño $O(h^3)$. Los productos con $\mu$ usan $\mu=O(h)$ y $\delta\mu=O(h^2)$; los del sitio elegido
usan $\hat c=O(h)$ y $\delta\hat c=O(h^2)$. En el punto~11, $\delta(6p^2+2p\alpha)=O(h^3)$; la suma
interior usa $\sum\tau|\mathsf G(x,e_0)|=O(1)$, y en los sitios retirados la masa es a lo sumo dos y
$\mathsf G(x,e_0)=O(h)$.

Para la traza se puede mantener $W$ exacta, con un coste $O(h^3)$ en cada tablero. Como $W$ anula el
calibre, la traza cuadrática es $\operatorname{tr}(G^0_{II}WG^0_{II}W)$. Las entradas de $W$ del
Paso~7 dan $\max|W|=O(h^2)$, $\|W\|_\infty=O(h)$ y, en los índices comunes, $\max|\delta W|=O(h^3)$ y
$\|\delta W\|_\infty=O(h^2)$; aquí se usan $\delta s=O(h^2)$, $\delta p=O(ph)$ y la fórmula de $Q_J$.
Escribamos $G^0_{II}=I+H_1$, con $\max|H_1|=O(h)$ y $\max|\delta H_1|=O(h^2)$. En
$\operatorname{tr}(W^2)$, cambiar $W$ cuesta $O(k\,h^2h^2)$ y perder una línea cuesta $O(h\,h^2)$. La
suma absoluta de $\operatorname{tr}(H_1W^2)$ es a lo sumo $k\max|H_1|\,\|W\|_\infty^2$, y la de
$\operatorname{tr}(H_1WH_1W)$ es a lo sumo $k^2\max|H_1|^2\|W\|_\infty^2$; cambiar un factor o fijar un
índice perdido mejora cada mayorante en un factor $h$. En total $\delta r_{e_0}=O(h^3)$, que es (a).
\end{proof}

\begin{corolario}[Incremento del resto de dos casillas]\label{cor:incremento-marginal}
En las hipótesis de la Proposición~\ref{prop:menor-lejano}\textup{(b)}, sea $R_\pi=\pi^T-\pi_1$ para el
par $\{e,f\}$. Si $c$ no está en ninguna de las filas ni de las columnas de $e,f$, entonces
\[
 |R_\pi(M_c)-R_\pi(M)|\le C\,\frac{p_ep_f}{k^3}.
\]
\end{corolario}
\begin{proof}
Por~\eqref{eq:psi-dos-casillas}, $R_\pi=p_ep_f\Psi$ con
$\Psi=(1+x_1)(1-g+t_1)(1+x_2+t_2)e^{\varrho_1+\varrho_2}-(1-g+x_1+x_2)$, donde $x_1=\hat c_e$,
$x_2=\hat c_f$, $g=\mathsf G(e,f)$, $t_1=p'_f/p_f-1+g$, $t_2=\hat c'_f-\hat c_f$, $\varrho_1=r_e(M)$ y
$\varrho_2=r_f(M_e)$. Aquí $x_i,g=O(h)$ y $t_i,\varrho_i=O(h^2)$. La parte de $\Psi$ lineal en
$(x_1,x_2,g)$ se anula, luego sus derivadas en $x_i,g$ son $O(h)$, mientras que las derivadas en
$t_i,\varrho_i$ son $1+O(h)$. Sea $\delta$ la variación de $M$ a $M_c$. Entonces
$\delta x_i,\delta g=O(h^2)$ por~\eqref{eq:variaciones-base}, y
$\delta\varrho_1,\delta\varrho_2=O(h^3)$ por la Proposición~\ref{prop:menor-lejano}\textup{(a)},
aplicada a $M$ y a $M_e$. Para $t_1,t_2$ hacen falta segundas diferencias. La matriz de Green del hijo
cumple $G'-G^0_{II}=(I+H_1)W(I+H_2)+C_q$, donde las entradas de $H_i$ son $O(h)$ con variaciones
$O(h^2)$, y $C_q$ se construye con $q=G^0_{IJ}v_J=O(h)$, $\delta q=O(h^2)$, $v_I^\top q=O(h)$ y el
divisor $2k-2$, de modo que $\delta C_q=O(h^3)$. Con las cotas de $W$ y $\delta W$ del Paso~3 resulta
$\delta(G'-G^0_{II})_{ij}=O(h^3)$. Para $X_x=(G^0a_x)^\top\eta^X-\varepsilon_x$ con
$\eta^X_i=\sum_{y\in i}\tau_y\varepsilon_y$, se tiene $\delta\Delta p=O(ph^2)$ por
$\Delta p=\tau s+O(ps^2)$ y~\eqref{eq:masas-hijo}, luego $\delta\Delta\varepsilon=O(h^3)$ y
$\delta\Delta\eta^X_i=O(h^3)$; en la contracción con $G^0a_x$, cuya norma $\ell^1$ es $O(1)$ y cuyas
entradas en las cuatro líneas retiradas son $O(h)$, esto da $\delta\Delta X=O(h^3)$ y
$\delta t_2=O(h^3)$. Por último,
$t_1=-(\varsigma_f/p_f)\mathsf G(e,f)+(\tau_f/p_f)a_f^\top\zeta+[u(t_fe^{s_f})-p_f-\tau_fs_f]/p_f$ tiene
variación $O(h^3)$ por el Paso~2. Por tanto $\delta\Psi=O(h^3)$. Como $\Psi=O(h^2)$ y
$\delta(p_ep_f)=O(p_ep_fh)$, se obtiene la tesis.
\end{proof}
